\subsection{LEXICOGRAPHIC PRODUCTS}

Let \((\mathbf{E}_t)_{i\in I}\) be a family of ordered sets, indexed by a well-ordered set \(I\) . Consider the product set \(\mathbf{E} = \prod_{i\in I}\mathbf{E}_t\) and the relation

`` \(x \in E\) and \(y \in E\) , and for the least index \(t \in I\) such that \(\mathrm{pr}_t x \neq \mathrm{pr}_t y\) , we have \(\mathrm{pr}_t x < \mathrm{pr}_t y\) '',

which we shall denote by \(\mathbf{R}\{x, y\}\) . It is evident that \(\mathbf{R}\{x, x\}\) is equivalent to \(x \in \mathbf{E}\) , that \(\mathbf{R}\{x, y\}\) implies \(\mathbf{R}\{x, x\}\) and \(\mathbf{R}\{y, y\}\) , and that \((\mathbf{R}\{x, y\}\) and \(\mathbf{R}\{y, x\})\) implies \(x = y\) . Also it is easily verified that \((\mathbf{R}\{x, y\}\) and \(\mathbf{R}\{y, z\})\) implies \(\mathbf{R}\{x, z\}\) (consider the least index \(t \in I\) for which at least two of the three elements \(\mathrm{pr}_t x\) , \(\mathrm{pr}_t y\) , \(\mathrm{pr}_t z\) are unequal); hence \(\mathbf{R}\{x, y\}\) is an order relation on the product set E. This relation and the ordering it defines are called the lexicographic order relation and the lexicographic ordering on E (induced by the given orderings on I and on the \(\mathbf{E}_t\) ); the set E with this ordering is called the lexicographic product of the family of ordered sets \((\mathbf{E}_t)_{t \in I}\) . If each \(\mathbf{E}_t\) is totally ordere'd the lexicographic product is also totally ordered.

